% =====================================================================
%  CHAPTER 5 -- IMPLEMENTATION AND TESTING
%
%  RUBRIC NOTE (R7): the top band requires implementation detail
%  sufficient "for the reader", plus testing performed "in very good
%  way" with enough test cases that "no further improvements are
%  required regarding the number and quality of test cases". In
%  practice: aim for a test-case table with real breadth -- happy paths,
%  boundary values, invalid input and security cases -- not four rows
%  that all pass.
% =====================================================================
\chapter{IMPLEMENTATION AND TESTING}
\label{ch:implementation-testing}

% TODO: one orientation paragraph.
[CHAPTER OPENING: state that this chapter describes how the design of
Chapter~\ref{ch:system-design} was implemented, and how the resulting
system was tested.]

% ---------------------------------------------------------------------
\section{IMPLEMENTATION}
\label{sec:implementation}
% ---------------------------------------------------------------------

\subsection{Development environment}
\label{subsec:dev-environment}

% Partly filled from the repository. The remaining bracketed values are
% facts about YOUR machines that no file records.
The system is composed of three independently runnable services. The
backend is a Laravel 12.57 application running on PHP 8.2. The frontend
is a React 18 and TypeScript application built with Vite. The NLP
service is a Python FastAPI application served by Uvicorn. Persistent
data is held in MySQL, and real-time messaging is carried over
WebSockets by Laravel Reverb. The full dependency set and versions are
listed in Table~\ref{tab:tools}.

% TODO: fill in the values below -- they cannot be read from the
% repository, and rubric R4 grades reproducibility.
The development environment used was
[OPERATING SYSTEM AND VERSION], with
[PYTHON VERSION] for the NLP service,
[NODE.JS VERSION] for the frontend toolchain, and
MySQL [VERSION] as the database server. Development hardware was
[BRIEF HARDWARE SPECIFICATION -- processor, RAM]. Visual Studio Code was
used as the primary IDE, and Git with GitHub for version control.

\subsection{Deployment architecture}
\label{subsec:deployment}

% TODO: your proposal names Vercel and Render. Confirm which service
% hosts which component, and whether the NLP service is deployed at all
% or only run locally -- that is a fair viva question.
[DEPLOYMENT: state which platform hosts the frontend, which hosts the
backend, where the MySQL database is hosted, and how the NLP service is
deployed. Note any environment variables or secrets that must be
configured, without listing their values.]

\subsection{Module implementation}
\label{subsec:module-implementation}

% TODO: one subsection or one paragraph per major module. For each
% module say what it does, which design component it realises, and the
% one genuinely non-obvious thing about how it was built.
[MODULE 1: name the module, state which component of
Section~\ref{sec:detailed-design} it implements, and describe how it
was built. Highlight the non-trivial part -- the piece that a reader
could not have guessed from the design chapter.]

[MODULE 2: as above.]

[MODULE 3: as above. Add one block per major module.]

\subsection{Key implementation details}
\label{subsec:key-details}

% TODO: cover the genuinely interesting engineering decisions -- the
% algorithm that needed optimising, the integration that was awkward,
% the security control you implemented, the migration you had to write.
[KEY DETAILS: describe the most significant implementation decisions
and any notable problems encountered and how they were resolved. If a
short code listing genuinely clarifies something, include it -- but
keep listings short and never paste whole files.]

% ---- Optional screenshot slots ---------------------------------------
% Screenshots of the working system count as supporting figures and are
% cheap marks under R10. Duplicate this block per screen.
\begin{figure}[H]
  \centering
  \fypdiagram[7cm]{[INSERT SCREENSHOT OF [SCREEN NAME] HERE]}
  \caption{[Describe what this screen does]}
  \label{fig:screenshot-1}
\end{figure}

\begin{figure}[H]
  \centering
  \fypdiagram[7cm]{[INSERT SCREENSHOT OF [SCREEN NAME] HERE]}
  \caption{[Describe what this screen does]}
  \label{fig:screenshot-2}
\end{figure}

% ---------------------------------------------------------------------
\section{TESTING}
\label{sec:testing}
% ---------------------------------------------------------------------

\subsection{Testing strategy}
\label{subsec:testing-strategy}

% TODO: name the levels of testing performed and the technique used at
% each level (black box, white box, boundary value analysis,
% equivalence partitioning).
[STRATEGY: state which levels of testing were carried out -- unit,
integration, system and user acceptance testing -- and the technique
applied at each level. State the tools used to run the tests.]

\subsection{Test cases}
\label{subsec:test-cases}

% ---------------------------------------------------------------------
%  TEST CASE TABLE
%  Uses longtable so it can run across several pages -- a table that
%  fills a page and stops is a common cause of lost marks here.
%  \small is used because six columns will not fit at 12pt on A4.
%  Only the first page of the table carries a caption, which is correct:
%  it means the List of Tables gets a single clean entry.
% ---------------------------------------------------------------------
{\small
\begin{longtable}{|l|p{2.6cm}|p{2.2cm}|p{2.4cm}|p{2.4cm}|l|}
  \caption{System test cases and their results}
  \label{tab:test-cases} \\
  \hline
  \textbf{Test ID} & \textbf{Description} & \textbf{Input} &
  \textbf{Expected Output} & \textbf{Actual Output} & \textbf{Status} \\
  \hline
  \endfirsthead

  % ---- header repeated at the top of every continuation page --------
  \hline
  \textbf{Test ID} & \textbf{Description} & \textbf{Input} &
  \textbf{Expected Output} & \textbf{Actual Output} & \textbf{Status} \\
  \hline
  \endhead

  \hline
  \endfoot

  % ---- TODO: replace these rows with your real test cases -----------
  TC-01 & [WHAT IS BEING TESTED] & [INPUT VALUES] &
         [EXPECTED RESULT] & [OBSERVED RESULT] & [Pass/Fail] \\
  \hline
  TC-02 & [WHAT IS BEING TESTED] & [INPUT VALUES] &
         [EXPECTED RESULT] & [OBSERVED RESULT] & [Pass/Fail] \\
  \hline
  TC-03 & [BOUNDARY VALUE CASE] & [INPUT VALUES] &
         [EXPECTED RESULT] & [OBSERVED RESULT] & [Pass/Fail] \\
  \hline
  TC-04 & [INVALID INPUT CASE] & [INPUT VALUES] &
         [EXPECTED RESULT] & [OBSERVED RESULT] & [Pass/Fail] \\
  \hline
  TC-05 & [SECURITY OR AUTHORISATION CASE] & [INPUT VALUES] &
         [EXPECTED RESULT] & [OBSERVED RESULT] & [Pass/Fail] \\
  \hline
  % TODO: keep adding rows. Cover every functional requirement in
  % Table 3.1, and include boundary, invalid-input and security cases --
  % a table of happy paths only will not reach the top band.
\end{longtable}
}

\subsection{Test summary}
\label{subsec:test-summary}

% TODO: summarise the outcome numerically.
[SUMMARY: how many test cases were executed, how many passed, how many
failed, and what was done about the failures. If any defect remains
open, say so honestly and state its severity -- an acknowledged known
issue reads far better than a silently omitted one.]
